\section{Flags + Limit High 字节（gran）}

字节 6 的 8 bit 被分成两半使用：

\begin{figure}[H]
  \centering
  % figures/ch03/fig06-gran-bit.tex — auto-extracted from chapter source
\begin{tikzpicture}
  \foreach \i/\label in {
    0/Lim$_{16}$, 1/Lim$_{17}$, 2/Lim$_{18}$, 3/Lim$_{19}$,
    4/AVL, 5/L, 6/D/B, 7/G} {
    \node[bitcell, minimum width=1.1cm, fill=white] (g\i) at (\i*1.1, 0) {\label};
    \node[font=\tiny\ttfamily, above=0.1cm of g\i] {bit \i};
  }
  \draw[decorate, decoration={brace, amplitude=5pt, mirror}]
    (0, -0.4) -- (3.3, -0.4) node[midway, below=5pt, font=\small] {limit[16:19]};
  \draw[decorate, decoration={brace, amplitude=5pt, mirror}]
    (3.3, -0.4) -- (8.7, -0.4) node[midway, below=5pt, font=\small] {flags};
\end{tikzpicture}

  \caption{gran 字节的 bit 布局}
\end{figure}

\begin{bitfield}
  \bit{7}{G (Granularity)：limit 的单位。
    G=0：limit 单位是字节（最大 $2^{20} - 1 = $ 1MiB）。
    G=1：limit 单位是 4KiB 页（最大 $2^{20} \times 4096 = $ 4GiB）。
    平坦模型 limit=\hex{FFFFF} + G=1 覆盖完整 4GiB 地址空间。}
  \bit{6}{D/B (Default Operation Size)：
    D=1：32 位段（默认操作数和地址大小是 32 位）。
    D=0：16 位段。我们在 32 位保护模式下必须设 D=1。}
  \bit{5}{L (Long mode)：64 位代码段标志。i386 下\textbf{必须为 0}。}
  \bit{4}{AVL (Available)：CPU 不使用，软件可自由使用。我们设为 0。}
  \bit{0--3}{limit 的高 4 位 [16:19]。}
\end{bitfield}

我们的 flags = \hex{C0}：

$$\hex{C0} = 1100\underbrace{0000}_{\text{limit\_hi}}$$

即 G=1, D/B=1, L=0, AVL=0。与 limit\_high = \hex{0F} 合并后 gran = \hex{CF}。

% ============================================================================

